% ============================================================
% GSPH-FC — Springer LNCS | WEA 2026
% Compilar: pdflatex -> bibtex -> pdflatex -> pdflatex
% ============================================================
\documentclass[runningheads]{llncs}
\usepackage[T1]{fontenc}
\usepackage[utf8]{inputenc}
\usepackage{graphicx}
\usepackage{amsmath,amssymb}
\usepackage{booktabs}
\usepackage{algorithm}
\usepackage[noend]{algpseudocode}
\usepackage{url}

\begin{document}

\title{GSPH-FC: A Geometric Space Partitioning Heuristic
with Frontier Connectivity for the Euclidean TSP}
\titlerunning{GSPH-FC: Geometric Space Partitioning for the TSP}

\author{%
    Martín Gómez Navarro\inst{1} \and
    Rodrigo Linfati\inst{2} \and
    Juan Carlos Figueroa-García\inst{3} \and
    Carlos Alberto Franco Franco\inst{3} \and
    Daniel Morillo-Torres\inst{4} \and
    Jairo R. Coronado-Hernandez\inst{5} \and
    Gustavo Esteban Gatica\inst{1}}
\authorrunning{M. Gómez Navarro et al.}
\institute{%
    Faculty of Engineering, Universidad Andres Bello, Santiago, Chile\\
    \email{m.gmeznavarro@uandresbello.edu}\\
    \email{ggatica@uandresbello.edu}
    \and
    Faculty of Engineering, Universidad del Bío-Bío,
    Concepción, Chile\\
    \email{rlinfati@ubiobio.cl}
    \and
    Faculty of Engineering, Universidad Distrital
    Francisco José de Caldas, Bogotá, Colombia\\
    \email{jcfigueroag@udistrital.edu.co}
    \email{cafrancof@udistrital.edu.co}
    \and
    Faculty of Engineering, Pontificia Universidad Javeriana de Cali,
    Cali, Colombia\\
    \email{daniel.morillo@javerianacali.edu.co}
    \and
    Faculty of Engineering, Universidad de la Costa,
    Barranquilla, Colombia\\
    \email{jcoronad18@cuc.edu.co}}

\maketitle

\begin{abstract}
We propose \textbf{GSPH-FC} (\textit{Geometric Space Partitioning with
Frontier Connectivity}), a geometric heuristic for the Euclidean
Traveling Salesman Problem (TSP). The method partitions the plane into
four quadrants, detects candidate cross-quadrant frontier pairs, and
constructs anchored intra-quadrant paths whose endpoints are fixed by
those frontier decisions. The resulting tour is assembled in geometric
order and refined by a strong time-driven polishing stage combining
candidate-set 2-opt, Or-opt with reversed insertion, and iterated local
search with double-bridge perturbations and localized reoptimization.
This yields an implementation that is mathematically consistent with the
frontier model and substantially stronger than a plain
endpoint-concatenation scheme. On the verified TSPLIB subset under
\texttt{EUC\_2D}, the final solver reaches the Best Known Solution (BKS)
exactly on many small and medium instances, including \texttt{pr124},
\texttt{pr144}, \texttt{pr152}, \texttt{ts225}, and \texttt{kroA200},
while remaining within 0.5\% of the BKS on harder cases such as
\texttt{a280}, \texttt{gil262}, and \texttt{lin318}. The method remains
geometrically interpretable and computationally competitive, while
delivering strong solution quality across the benchmark subset studied
here.
\keywords{2-opt \and Frontier connectivity \and Geometric decomposition
\and Space partitioning \and TSPLIB benchmarks \and Traveling
Salesman Problem}
\end{abstract}

% ================================================================
\section{Introduction}
% ================================================================

The Euclidean \textit{Traveling Salesman Problem} (TSP) asks for a
minimum-length Hamiltonian cycle through $n$ planar points and remains
one of the most classical testbeds for heuristic design in
combinatorial optimization~\cite{applegate2006tsp,laporte1992tsp}.
Although exact methods solve many benchmark instances to optimality,
heuristics remain relevant whenever one seeks interpretable, scalable,
and lightweight solvers with explicit design structure rather than a
large amount of hidden machinery~\cite{garey1979computers,mardones2023}.

The central hypothesis behind \textbf{GSPH-FC} (\textit{Geometric Space
Partitioning with Frontier Connectivity}) is that geometric structure
should guide not only the \emph{placement} of subproblems but also the
\emph{assembly} of their local solutions. In the final method,
frontier connectivity is treated as a first-class construction
principle: frontier pairs are selected before local routes are built and
used as mandatory entry/exit anchors for the intra-quadrant paths.
Hence, the global assembly phase is directly aligned with the geometric
frontier model rather than with a generic endpoint-concatenation rule.

The final solver has three layers. First, it generates a small portfolio
of geometric partitions around the bounding-box midpoint. Second, for
each partition it computes candidate frontier pairs and builds anchored
open paths inside each quadrant. Third, it applies a strong global
polishing phase based on time-driven iterated local search (ILS),
candidate-set 2-opt, and Or-opt with reversed insertion. This retains
the geometric motivation of the method while making the resulting tour
quality substantially stronger.

\paragraph{Contributions.}
\begin{enumerate}
  \item \textbf{Frontier-guided anchored assembly:} a geometric TSP
    heuristic in which frontier pairs are selected before local route
    construction and define the entry/exit anchors of each quadrant.
  \item \textbf{Strong unified polishing:} a time-driven ILS with
    localized reoptimization of perturbed regions, candidate-set 2-opt,
    and Or-opt with reversed insertion.
  \item \textbf{Updated empirical evidence:} extensive results on a
    verified TSPLIB subset showing many exact BKS matches on small and
    medium instances and very small residual gaps on harder medium and
    large cases.
\end{enumerate}

The remainder is organized as follows. Section~\ref{sec:related}
reviews the relevant background. Section~\ref{sec:formulation}
presents the geometric formulation. Section~\ref{sec:algorithm}
summarizes the final algorithm. Section~\ref{sec:experiments} reports
the updated experimental results. Section~\ref{sec:analysis} discusses
strengths and limitations, and Section~\ref{sec:conclusions}
concludes.

% ================================================================
\section{Related Work}
\label{sec:related}
% ================================================================

GSPH-FC is built on ideas from geometric TSP heuristics, local search,
and trajectory metaheuristics. Nearest Neighbour (NN) is the classical
constructive baseline: it runs in $O(n^2)$ and is often used to seed
improvement heuristics, but by itself it is usually far from
optimal~\cite{kizilatecs2013nearest}. The 2-opt heuristic improves a
seed tour by reversing segments that shorten total
length~\cite{croes1958method}; together with candidate sets and
Don't-Look Bits, it remains one of the most important building blocks
for Euclidean local search~\cite{johnson1997tsp,bentley1992fast}. The
Lin--Kernighan family and its descendants such as LKH extend this idea
to stronger variable-depth neighborhoods and are among the strongest
practical TSP heuristics available~\cite{lin1973effective,helsgaun2000effective,helsgaun2017extension}.

At a higher level, many metaheuristics alternate perturbation and local
improvement to escape local optima. Iterated Local Search (ILS) is
especially relevant because it couples strong local search with
lightweight diversification through perturbation and restart
strategies~\cite{lourencco2003ils}. On TSP, double-bridge kicks plus
2-opt/3-opt style reoptimization are a particularly successful
combination~\cite{martin1991large}. GSPH-FC follows this line, but adds
an explicit geometric decomposition layer and makes frontier
connectivity part of the local route construction itself.

Our claim is therefore not that NN, 2-opt, Or-opt, or ILS are novel in
isolation. Rather, the contribution lies in the way geometric
partitioning, frontier-guided anchoring, and strong localized ILS are
combined into a coherent geometric solver.

\begin{table}
\caption{Main TSP method families and their qualitative tradeoffs.}
\label{tab:comparison}
\centering\small
\begin{tabular}{@{}lllll@{}}
\toprule
\textbf{Method} & \textbf{Type} & \textbf{Complexity} & \textbf{Scalability} & \textbf{Quality} \\
\midrule
Brute Force         & Exact         & $O(n!)$      & Very low  & Optimal \\
Dynamic Programming & Exact         & $O(n^2 2^n)$ & Low       & Optimal \\
Nearest Neighbour   & Constructive  & $O(n^2)$     & High      & Low \\
2-opt               & Improvement   & $O(n^2)$/iter& Medium    & Medium \\
Lin--Kernighan/LKH  & Improvement   & Variable     & High      & Very high \\
ILS-based heuristics& Metaheuristic & Variable     & High      & High \\
\textbf{GSPH-FC}   & Geometric heuristic & $O(n^2)$ & High & High \\
\bottomrule
\end{tabular}
\end{table}

% ================================================================
\section{Mathematical Formulation}
\label{sec:formulation}
% ================================================================

Let $G=(V,E)$ be a complete undirected graph where
$V=\{v_1,\dots,v_n\}\subset\mathbb{R}^2$ with coordinates
$(x_i,y_i)$. Edge costs follow the TSPLIB \texttt{EUC\_2D}
convention~\cite{reinelt1991tsplib}:
\begin{equation}
 d(v_i,v_j)=\left\lfloor\sqrt{(x_i-x_j)^2+(y_i-y_j)^2}+0.5\right\rfloor.
 \label{eq:dist}
\end{equation}
The TSP seeks a Hamiltonian cycle
$\pi=(v_{\pi(1)},\dots,v_{\pi(n)},v_{\pi(1)})$ minimizing
\begin{equation}
 \min_{\pi}\;L(\pi)=\sum_{i=1}^{n}
 d\!\left(v_{\pi(i)},v_{\pi(i\bmod n+1)}\right).
 \label{eq:tsp}
\end{equation}

\paragraph{Scope and assumptions.}
We restrict attention to planar Euclidean instances under the
\texttt{EUC\_2D} rounding convention. All experimental GAP values are
computed against BKS values under the same metric.

\subsection{Geometric Decomposition}

Let the bounding-box midpoint be
\begin{equation}
 x_m=\tfrac{x_{\min}+x_{\max}}{2}, \qquad
 y_m=\tfrac{y_{\min}+y_{\max}}{2}.
\end{equation}
The instance is partitioned into four quadrants:
\begin{align}
Q_1&=\{v\in V: x_v\le x_m,\; y_v>y_m\}  &&\text{(NW)},\\
Q_2&=\{v\in V: x_v>x_m,\; y_v>y_m\}    &&\text{(NE)},\\
Q_3&=\{v\in V: x_v\le x_m,\; y_v\le y_m\}&&\text{(SW)},\\
Q_4&=\{v\in V: x_v>x_m,\; y_v\le y_m\} &&\text{(SE)}.
\end{align}
In the final solver, this basic partition can be slightly shifted to
form a small portfolio of candidate partitions, but each candidate
remains a four-way axis-aligned decomposition.

\subsection{Frontier Connectivity}

For adjacent quadrants $Q_a,Q_b$, define a frontier band around the
shared boundary. For a vertical frontier at $x_m$:
\begin{equation}
F_a(\varepsilon)=\{u\in Q_a:|x_u-x_m|<\varepsilon\}, \qquad
F_b(\varepsilon)=\{v\in Q_b:|x_v-x_m|<\varepsilon\}.
\end{equation}
The best frontier pair is then
\begin{equation}
(u^*,v^*)=\arg\min_{u\in F_a,\,v\in F_b} d(u,v).
\label{eq:frontier}
\end{equation}
If a frontier band is empty, $\varepsilon$ is expanded adaptively.
Crucially, these frontier pairs are chosen \emph{before} local route
construction so that they define the actual anchors of the local paths.

\subsection{Anchored Intra-Quadrant Construction}

Given a quadrant $Q_k$ and two anchors $s_k,t_k\in Q_k$, the local
subroute is built as an \emph{open anchored path} whose endpoints are
fixed at $(s_k,t_k)$. The interior is completed by minimum-cost
insertion and refined by 2-opt on the open path. Thus the local
construction is conditioned by the inter-quadrant connectivity choices,
aligning the mathematical frontier model with an implementable route
construction scheme.

\subsection{Global Assembly and Strong Polishing}

Quadrant paths are assembled in the geometric order
$Q_1\!\to\!Q_2\!\to\!Q_4\!\to\!Q_3$, producing one or more plausible
seed tours. These are filtered cheaply and only the best candidates are
passed to a strong global polishing stage.

The final polishing combines: (i) candidate-set 2-opt with
Don't-Look Bits, (ii) Or-opt relocation with optional reversed
insertion, and (iii) time-driven ILS with double-bridge perturbations.
In the current solver, ILS is not capped by a fixed number of
perturbations; instead it continues while budget remains. After each
perturbation, only the perturbed region and its candidate neighbors are
reactivated for local reoptimization, making the intensification
substantially stronger per unit of time.

\subsection{Complexity}

The geometric partition is $O(n)$, candidate sets are built in
$O(n\log n)$ or $O(nk)$ depending on the backend, and anchored local
construction is quadratic in the quadrant size in the worst case. With
sparse candidate sets, the dominant practical cost is the strong global
polishing, which remains in the quadratic regime typical of strong
2-opt-based Euclidean TSP heuristics. The method therefore preserves a
manageable asymptotic profile while investing the computational budget
where it is most effective: the final global intensification.

% ================================================================
\section{Algorithm Overview}
\label{sec:algorithm}
% ================================================================

This section summarizes the final workflow of GSPH-FC. The first
algorithm describes how candidate partitions and frontier choices are
combined; the second shows how each quadrant is built under fixed
anchors; the third summarizes the global intensification stage.

\paragraph{Top-level workflow.}
The main procedure generates a small set of geometrically plausible seed
tours and spends the main computational effort on polishing only the
best candidates.

\begin{algorithm}
\caption{GSPH-FC --- top-level procedure}
\label{alg:gsph}
\begin{algorithmic}[1]
\Require $V\subset\mathbb{R}^2$, frontier tolerance $\varepsilon$
\Ensure Tour $\pi$, length $L$
\State Generate a small portfolio of geometric partitions
\ForAll{partition candidate $(x_m,y_m)$}
  \State Assign nodes to $Q_1,Q_2,Q_3,Q_4$
  \State Compute frontier candidates on
         $(Q_1,Q_2),(Q_2,Q_4),(Q_4,Q_3),(Q_3,Q_1)$
  \ForAll{selected frontier combinations}
    \ForAll{$Q_k\in\{Q_1,Q_2,Q_4,Q_3\}$}
      \State $R_k\gets\Call{AnchoredPath}{Q_k,\text{entry}(Q_k),\text{exit}(Q_k)}$
    \EndFor
    \State Assemble the seed tour $\pi_0$ in geometric order
    \State Apply light polishing and keep only the best seed tours
  \EndFor
\EndFor
\State Apply \Call{StrongPolish}{} to the finalists
\State \Return the best polished tour
\end{algorithmic}
\end{algorithm}

\paragraph{Anchored local construction.}
This routine is the key mechanism that makes frontier connectivity
govern the intra-quadrant paths rather than being used only after local
construction is finished.

\begin{algorithm}
\caption{\textsc{AnchoredPath} --- anchored intra-quadrant construction}
\label{alg:anchored}
\begin{algorithmic}[1]
\Require Quadrant point set $Q$, entry anchor $s$, exit anchor $t$
\Ensure Open path $P$
\If{$|Q|\le 1$}
  \State \Return $Q$
\EndIf
\If{$s$ or $t$ is invalid}
  \State Build a greedy NN path and refine it with 2-opt
  \State \Return the improved path
\EndIf
\State Initialize $P=[s,t]$
\While{unvisited nodes remain in $Q$}
  \State Insert the node and position with minimum insertion cost
\EndWhile
\State Refine the open path with 2-opt while keeping the anchors fixed
\State \Return $P$
\end{algorithmic}
\end{algorithm}

\paragraph{Strong global polishing.}
The final stage intensifies the best seed tours with a time-driven ILS
that repeatedly perturbs and reoptimizes only the affected region.

\begin{algorithm}
\caption{\textsc{StrongPolish} --- time-driven global intensification}
\label{alg:polish}
\begin{algorithmic}[1]
\Require Seed tour $\pi_0$, budget $T_{\text{bgt}}(n)$
\Ensure Strongly polished tour
\State Build distance matrix $D$ and candidate sets $C$
\State Reach an initial local optimum with 2-opt and Or-opt
\State $\pi^*\gets$ incumbent best tour
\While{$\Call{Now}{}<T_{\text{bgt}}(n)$}
  \State Generate a double-bridge perturbation of the current tour
  \State Reactivate only perturbed nodes and their candidate neighbors
  \State Reapply localized 2-opt and Or-opt
  \If{the new tour improves $\pi^*$}
    \State update the incumbent and intensify from it
  \Else
    \State continue from the current incumbent or restart after stagnation
  \EndIf
\EndWhile
\State \Return $\pi^*$
\end{algorithmic}
\end{algorithm}

% ================================================================
\section{Experimental Results}
\label{sec:experiments}
% ================================================================

\paragraph{Setup.}
All results below use the same final GSPH-FC configuration with
anchored frontier construction and strong polishing. For profile
experiments we report the best tour over 5 runs under the
\texttt{EUC\_2D} metric. GAP$(\%)=100\times
(\text{GSPH}-\text{BKS})/\text{BKS}$. Instances without an agreed BKS
are marked \textbf{N/A}. We only report results that were verified with
the final solver.

\begin{table}
\caption{Small instances ($n\le 120$) under the final GSPH-FC solver.
Bold GAP indicates an exact BKS match.}
\label{tab:small}
\centering\scriptsize
\begin{tabular}{@{}lrrrrr@{}}
\toprule
\textbf{Instance} & $n$ & \textbf{GSPH} & \textbf{BKS} & \textbf{GAP (\%)} & \textbf{Time (s)}\\
\midrule
berlin52 &  52 &   7{,}542 &   7{,}542 & \textbf{0.000} & 10.092 \\
 eil51    &  51 &       426 &       426 & \textbf{0.000} & 10.046 \\
 eil76    &  76 &       538 &       538 & \textbf{0.000} & 10.136 \\
 eil101   & 101 &       630 &       629 & 0.159 & 10.299 \\
 eil10    &  10 &       159 &        --- & \textbf{N/A} & 10.002 \\
 eil20    &  20 &        --- &        --- & \textbf{N/A} & --- \\
 rd100    & 100 &   7{,}910 &   7{,}910 & \textbf{0.000} & 5.263 \\
 rat99    &  99 &   1{,}212 &   1{,}211 & 0.083 & 10.268 \\
 st70     &  70 &       675 &       675 & \textbf{0.000} & 10.093 \\
 pr76     &  76 & 108{,}159 & 108{,}159 & \textbf{0.000} & 10.145 \\
 kroA100  & 100 &  21{,}282 &  21{,}282 & \textbf{0.000} & 5.223 \\
 kroB100  & 100 &  22{,}141 &  22{,}141 & \textbf{0.000} & 5.234 \\
 kroC100  & 100 &  20{,}749 &  20{,}749 & \textbf{0.000} & 5.267 \\
 kroD100  & 100 &  21{,}294 &  21{,}294 & \textbf{0.000} & 10.276 \\
 kroE100  & 100 &  22{,}106 &  22{,}068 & 0.172 & 10.271 \\
 lin105   & 105 &  14{,}379 &  14{,}379 & \textbf{0.000} & 10.424 \\
 pr107    & 107 &  44{,}387 &  44{,}303 & 0.190 & 10.306 \\
\midrule
\multicolumn{4}{l}{\textit{Avg. GAP (excl. N/A)}} & \textbf{0.040} & 8.890 \\
\bottomrule
\end{tabular}
\end{table}

\begin{table}
\caption{Medium instances ($121\le n\le 300$) under the final GSPH-FC solver.
Bold GAP indicates an exact BKS match.}
\label{tab:medium}
\centering\scriptsize
\begin{tabular}{@{}lrrrrr@{}}
\toprule
\textbf{Instance} & $n$ & \textbf{GSPH} & \textbf{BKS} & \textbf{GAP (\%)} & \textbf{Time (s)}\\
\midrule
bier127  & 127 & 118{,}282 & 118{,}282 & \textbf{0.000} & 5.867 \\
 ch130    & 130 &   6{,}110 &   6{,}110 & \textbf{0.000} & 10.564 \\
 ch150    & 150 &   6{,}528 &   6{,}528 & \textbf{0.000} & 10.844 \\
 pr124    & 124 &  59{,}030 &  59{,}030 & \textbf{0.000} & 10.424 \\
 pr136    & 136 &  96{,}785 &  96{,}772 & 0.013 & 10.545 \\
 pr144    & 144 &  58{,}537 &  58{,}537 & \textbf{0.000} & 10.652 \\
 pr152    & 152 &  73{,}682 &  73{,}682 & \textbf{0.000} & 24.766 \\
 ts225    & 225 & 126{,}643 & 126{,}643 & \textbf{0.000} & 25.494 \\
 tsp225   & 225 &   3{,}920 &   3{,}916 & 0.102 & 25.632 \\
 a280     & 280 &   2{,}589 &   2{,}579 & 0.388 & 15.238 \\
 kroA150  & 150 &  26{,}524 &  26{,}524 & \textbf{0.000} & 5.644 \\
 kroB150  & 150 &  26{,}132 &  26{,}130 & 0.008 & 10.763 \\
 pr226    & 226 &  80{,}369 &  80{,}369 & \textbf{0.000} & 25.282 \\
 pr264    & 264 &  49{,}378 &  49{,}135 & 0.495 & 27.363 \\
 pr299    & 299 &  48{,}224 &  48{,}191 & 0.068 & 27.550 \\
 kroA200  & 200 &  29{,}368 &  29{,}368 & \textbf{0.000} & 25.005 \\
 kroB200  & 200 &  29{,}438 &  29{,}437 & 0.003 & 25.049 \\
 gil262   & 262 &   2{,}387 &   2{,}378 & 0.378 & 26.267 \\
 d198     & 198 &  15{,}793 &  15{,}780 & 0.082 & 14.885 \\
 u159     & 159 &  42{,}080 &  42{,}080 & \textbf{0.000} & 12.678 \\
\midrule
\multicolumn{4}{l}{\textit{Avg. GAP}} & \textbf{0.077} & 17.526 \\
\bottomrule
\end{tabular}
\end{table}

\begin{table}
\caption{Representative harder instances under the final solver.}
\label{tab:hard}
\centering\small
\begin{tabular}{@{}lrrrrr@{}}
\toprule
\textbf{Instance} & $n$ & \textbf{GSPH} & \textbf{BKS} & \textbf{GAP (\%)} & \textbf{Time (s)}\\
\midrule
pr136   & 136 &  96{,}785 &  96{,}772 & 0.013 & 10.545 \\
 a280    & 280 &   2{,}589 &   2{,}579 & 0.388 & 15.238 \\
 gil262  & 262 &   2{,}387 &   2{,}378 & 0.378 & 26.267 \\
 lin318  & 318 &  42{,}208 &  42{,}029 & 0.426 & 27.328 \\
\bottomrule
\end{tabular}
\end{table}

\begin{table}
\caption{Compact summary by size range over the verified subset reported here.}
\label{tab:summary}
\centering\small
\begin{tabular}{@{}lrrrr@{}}
\toprule
\textbf{Range} & \textbf{Cases} & \textbf{Avg. GAP (\%)} & \textbf{Max GAP (\%)} & \textbf{Avg. time (s)}\\
\midrule
$n\le 120$        & 15 & 0.040 & 0.190 & 8.890 \\
$121\le n\le 300$ & 20 & 0.077 & 0.495 & 17.526 \\
$301\le n\le 800$ &  1 & 0.426 & 0.426 & 27.328 \\
\bottomrule
\end{tabular}
\end{table}

The final solver is strongest on small and medium Euclidean TSPLIB
instances, where it frequently matches the BKS exactly. In
Table~\ref{tab:small}, exact matches dominate the verified small set;
only a few cases such as \texttt{eil101}, \texttt{rat99},
\texttt{kroE100}, and \texttt{pr107} remain slightly above the BKS.
Table~\ref{tab:medium} shows that the same behavior extends into the
medium range: exact matches occur for \texttt{bier127}, \texttt{ch130},
\texttt{ch150}, \texttt{pr124}, \texttt{pr144}, \texttt{pr152},
\texttt{ts225}, \texttt{kroA150}, \texttt{pr226}, \texttt{kroA200}, and
\texttt{u159}, while most remaining cases still stay below 0.5\% GAP.

The harder representative cases in Table~\ref{tab:hard} show that the
method remains highly competitive beyond the easiest benchmark regime.
In particular, \texttt{a280} reaches 0.388\% GAP, \texttt{gil262}
reaches 0.378\%, and \texttt{lin318} reaches 0.426\% on the verified
profiled runs. These cases are important because they show that the
frontier-guided anchored construction still yields strong seed tours on
more difficult instances, and that the time-driven ILS closes most of
the remaining gap. Figures~\ref{fig:lin318-route} and
\ref{fig:krob100-route} provide two qualitative views of the final
solver: one harder large instance and one medium-scale benchmark where
the geometric structure of the resulting route is easy to inspect.

\begin{figure}
\centering
\includegraphics[width=0.92\textwidth]{Resultados/LIN318.png}
\caption{Representative large-instance result on \texttt{lin318}.
The figure illustrates the structured global tour produced by the
final GSPH-FC solver on a harder Euclidean benchmark.}
\label{fig:lin318-route}
\end{figure}

\begin{figure}
\centering
\includegraphics[width=0.92\textwidth]{Resultados/KROB100.png}
\caption{Representative medium-scale result on \texttt{kroB100}.
This instance illustrates the high-quality final tour structure
obtained by GSPH-FC on a standard 100-city benchmark.}
\label{fig:krob100-route}
\end{figure}

Taken together, the results indicate that the geometric decomposition is
most effective when paired with strong global intensification. The
construction phase yields structured low-crossing seed tours, while the
time-driven ILS and localized reoptimization are what push the method
from a reasonable heuristic to a solver that often reaches the BKS
exactly on small and medium instances and stays very close on harder
ones.

% ================================================================
\section{Discussion}
\label{sec:analysis}
% ================================================================

The strongest regime of GSPH-FC is the small-to-medium Euclidean TSP,
where frontier-guided anchoring and localized ILS often suffice to
reach the exact BKS. In this range, the geometric partition provides
useful structure without over-constraining the search, and the global
polishing stage can intensify efficiently on a well-shaped seed tour.

The main limitation is that axis-aligned partitioning may still be
suboptimal for elongated or diagonally clustered instances. In such
cases, even multiple frontier candidates can leave a few long global
links that require stronger post-assembly search to remove. The method
is also metric-specific: all reported results assume \texttt{EUC\_2D},
and historical BKS values under a different convention must be handled
carefully.

Promising next steps include adaptive partition orientations,
frontier-focused higher-order local search (e.g., selective 3-opt or
LK around difficult links), and stronger candidate structures for the
largest instances. These directions preserve the geometric
interpretability of GSPH-FC while targeting the remaining gap on harder
large-scale cases.

% ================================================================
\section{Conclusions}
\label{sec:conclusions}
% ================================================================

We presented the final version of \textbf{GSPH-FC}, a Euclidean TSP
heuristic whose defining feature is to make frontier connectivity a
first-class construction principle. Frontier pairs are selected before
local route construction, used as explicit entry/exit anchors for the
quadrants, and followed by a strong time-driven polishing stage based
on localized 2-opt, reversed Or-opt, and ILS with double-bridge
perturbations. This yields a solver that is both mathematically
coherent and empirically strong within the verified TSPLIB subset.

Across the verified TSPLIB subset reported in this paper, the final
solver reaches the exact BKS on many small and medium instances and
remains very close on harder cases such as \texttt{a280},
\texttt{gil262}, and \texttt{lin318}. The resulting method remains
conceptually simple and geometrically interpretable, while delivering
highly competitive solution quality. Future work will focus on adaptive
partition orientations and selective higher-order local search near
difficult frontier regions.

% BLIND REVIEW: acknowledgements hidden
%
% RESTORE BEFORE FINAL SUBMISSION:
%\begin{credits}
%\subsubsection{\ackname}
%The authors thank the faculty of the Heuristics course at
%Universidad Andres Bello for their guidance and material support
%during the development of this work.
%\end{credits}

\bibliographystyle{splncs04}
\bibliography{bibliografia}

\end{document}
